\section{Encoding, computation, and the neighboring problems}\label{sec:encoding49}
The source of an error certificate must be distinguished from the cost
of verifying a supplied realization. The new alternatives require the
complete numerical tensor. Its size is $n d\prod_t|\calA_t|$; it need
not be polynomial in a horizon described without its words. The
all-word Gram recurrence below, by contrast, accepts an explicitly
given orthogonal target and a proposed machine and need not expand
the word table. It decides exactness of that candidate, not whether
an optimum over candidates has been found.

\begin{proposition}[An inherited hardness boundary in a binary interface]
\label{prop:hard49}
For complete rational antisymmetric tables with one command epoch,
one selected query, and a variable seed and command alphabet,
deciding $\mathcal E_2(f)\le\varepsilon$ for rational $\varepsilon$
is NP-complete. Here $\varepsilon$ is binary-TV tolerance and $2\varepsilon$ is mean tolerance. This is a direct transfer of the rank-one matrix
infinity-norm decision theorem of Gillis--Shitov \cite[Theorem 3]{GS19},
not a new hardness theorem for a fixed-dimensional orthogonal alphabet.
\end{proposition}
\begin{proof}
Given the matrix instance $(M,k)$ in that theorem, let
$B=\max(1,\|M\|_\infty+k)$ and prescribe $f_{i,a}=M_{ia}/B$
and tolerance $\varepsilon=k/(2B)$. If $X=uv^T$ approximates $M$
within $k$, then $X/B$ has maximum entry magnitude at most one.
For a rank-one matrix this maximum is the product of the two factor
infinity norms. Rescale one factor to infinity norm one and the
other to at most one; the zero matrix is treated separately.
The scalar normal form gives the required two-state realization.
Conversely its two bounded factors give a rank-one approximation to
$M$. All input numbers and tables in the reduction have polynomial
binary length.

For membership, strip threshold-inactive rows and columns as in
Lemma~\ref{lem:support49} and guess their nonzero factor signs.
Writing the remaining positive magnitudes as $h_i\beta_a$ and
$t_i=1/h_i\ge1$, the constraints become
\[
 (\sigma_{ia}f_{ia}-2\varepsilon)t_i\le\beta_a
 \le(\sigma_{ia}f_{ia}+2\varepsilon)t_i,
 \qquad t_i\ge1,\quad0\le\beta_a\le1.
\]
This is a rational linear system. On essential columns the active
constraints force $\beta_a>0$ automatically. A nonempty rational
polyhedron of this form has a rational point with polynomial bit
length, obtained from an appropriate basic system by Cramer's rule: after clearing denominators, $q$ variables and coefficient bit length $L$ give determinant bit length $O(q(L+\log(q+1)))$. The feasible set is pointed because $t_i\ge1$ and $0\le\beta_a\le1$, so a basic feasible point exists. This is the standard rational-polyhedron argument; compare \cite[Chapter 10]{Sch86}. Its reciprocals
and the guessed signs give a polynomially checkable rational machine.
The single query factor is absorbed into the command factors.
Thus this restricted decision problem is in NP.
\end{proof}

This result explains why enumerating all sign chambers must not be
called a general polynomial-time algorithm. Even the matrix problem
already has a substantial optimization literature. Gillis--Shitov,
Lemmas 1--2, remove threshold-inactive coordinates and reduce a
fixed-sign matrix instance to linear inequalities using reciprocal
variables. Morozov--Smirnov--Zamarashkin analyze alternating
minimization and sign-transition graphs for global rank-one Chebyshev
approximation \cite[Theorem 5.3]{MSZ23}. Neither fixed-sign feasibility
nor global matrix optimization is introduced here as a discovery.
For the multi-epoch tensor, logarithmic magnitudes linearize its
higher-order incidence constraints. The construction is the familiar
monomial-to-linear conversion of geometric programming; the same
reference already describes threshold bisection for relative fitting
error \cite[\S2 and \S8]{BKVH07}. Our objective remains \emph{absolute}
mean error, encoded by its exact asymmetric lower and upper endpoints,
not the log-Chebyshev objective obtained by fitting logarithms of the
data.

The difference asserted in this revision is the complete finite
alternative for the controlled bounded tensor, its zero-boundary
handling, its multiplicative infeasibility witnesses against all
machines in the two-state class, and the reconstruction of legal
stochastic rows. The exact cubic example is an application in which
sign consistency alone is insufficient. The existence of such
certificates is proved from classical linear alternatives, not from
an asserted novel version of Farkas' lemma. The multiple-cut
compatibility tradeoff of Theorem~\ref{thm:frontier47} and the original
all-hidden-state lower bounds are inherited and retained.

\begin{center}
\small
\begin{tabular}{p{.23\textwidth}p{.30\textwidth}p{.36\textwidth}}
\toprule
Question & Relevant antecedent & Statement used here\\\midrule
Candidate equality & Weighted/probabilistic equivalence
\cite{Tzeng92,BPP15} & A source-bound finite-clock Gram test, not an optimum certificate.\\[3pt]
Positive realization & Invariant cones and controlled suffixes
\cite{BF04,GTF33} & Compatible positive Hankel factors retained with all suffixes.\\[3pt]
Fixed-sign approximation & Matrix Chebyshev feasibility
\cite{GS19} and monomial logarithms \cite{BKVH07} & Exact multi-mode bounds, zero support, sign chambers and multiplicative alternatives.\\[3pt]
Global two-state error & Scalar reduction and sign obstruction in revision 47
\cite{GTF47} & Complete magnitude certificates and an exact cubic numerical optimum.\\[3pt]
Arithmetic width & Prior packet and Diophantine proofs
\cite{GTF35,GTF42,GTF43} & Retained, not assigned new exponents by the present certificate theorem.\\
\bottomrule
\end{tabular}
\end{center}

The lower optimization concerns two available labels at each selected
command boundary. This two-state optimizer does not classify general three-state or larger hidden realizations. The distinct rank-tight geometric results of Section~\ref{sec:simplex50} do apply beyond two states under their stated spanning and orthogonality hypotheses. Common positive Hankel factorization remains the
appropriate general model; the one-dimensional odd coordinate is
special to two states. Nor does the finite algorithm bound ordinary
uniform workspace: it charges neither stored row descriptions nor
arithmetic, and it does not turn an algebraic optimum into exact
finite-fair-bit sampling. The separate compiler in the retained
material states its own assumptions and space accounting.
